\section{Three comparison scales and weak convergence}
\label{sec:three-panel-weak-comparison}

The comparison estimate is initially available at grid points whose coordinate
projections are distinct. A test function may cross the excluded hyperplanes.
We use three scales to pass from the former estimate to the latter test:
one for the numerical comparison, one for a bound everywhere on a compact
noncoincident set, and one for a derivative bound away from thin strips.
Their different costs are essential to the argument.

\subsection{The geometric and analytic estimates}

Fix an arity $m\ge1$ and use the Euclidean norm on
$E_4^m=(\mathbb R^4)^m$, retaining the maximum norm for pair separation. For $R>0$ set
\[
 K_R=\{x:\|x\|_2\le R,\quad
               \|x_i-x_j\|_\infty\ge R^{-1}\ (i\ne j)\}.
\]
This is a compact subset of the noncoincident configuration space. Every
compact subset of that space is contained in some $K_R$. Define
\[
 \Omega_\sigma=\{x: |x_{i,\nu}-x_{j,\nu}|>8\sigma
                  \text{ for every }i\ne j,\ 1\le\nu\le4\}.
\]
Thus the complement of $\Omega_\sigma$ is a union of coordinate strips,
not a neighborhood of just the actual coincidence set.

Let $S$ be the smooth interacting table already extracted along the ordered
mesh and volume filter $\mathcal F$, and let $Q(x)$ be the free correlation
for a fixed word of labels. Fix a comparison index $k$. Its physical radius
is $r_k>0$; its normalized grid spacing is $h_k>0$; its strip width is
$\sigma_k>0$. The calibrated insertion maps may still depend on the finite
cutoff pair. Suppose they are eventually bounded in their finite-dimensional
operator space. On any convergent subfilter of these maps write $Y$ for the
limiting tuple and
\[
                 U_{k,Y}(x)=S(r_kx;Y).
\]
All the following estimates are to hold for every such $Y$, with constants
independent of that subfilter. The three required estimates are
\begin{align}
 |U_{k,Y}(y)-Q(y)|&\le e_k
 &&\text{at generic grid points }y\in K_{2R},\label{eq:three-grid}\\
 |U_{k,Y}(x)|&\le A_k^m
 &&(x\in K_R),\label{eq:three-bulk}\\
 \|D U_{k,Y}(x)\|&\le 4m B_k^m/\sigma_k
 &&\bigl(|x_{i,\nu}-x_{j,\nu}|>4\sigma_k\bigr).
 \label{eq:three-slope}
\end{align}
The last estimate is needed only in the noncoincident domain. The quantities
$e_k,A_k,B_k$ are nonnegative. Assume
\begin{gather}
 e_k\longrightarrow0,\qquad \sigma_k\longrightarrow0,\qquad
                 4m h_k\le\sigma_k\quad\hbox{eventually},\label{eq:three-room}\\
 A_k^m\sigma_k\longrightarrow0,\qquad
                  B_k^m h_k/\sigma_k\longrightarrow0.
                  \label{eq:three-balances}
\end{gather}
We first prove the weak comparison from these estimates, then explain the
three scale choices that give exactly these powers.

Choose the coordinatewise lattice address $y$ of $x$, so
$\|y-x\|\le4m h_k$. If $x\in\Omega_{\sigma_k}$, the closed ball
of radius $\sigma_k$ about $x$ has every coordinate difference larger
than $4\sigma_k$: moving its two endpoints changes a difference by at
most $2\sigma_k$. In particular the segment from $x$ to $y$ remains
in the region of \eqref{eq:three-slope}. The fundamental theorem of
calculus on that segment gives
\[
 |U_{k,Y}(x)-U_{k,Y}(y)|
          \le16m^2 B_k^m h_k/\sigma_k.
\]
For all sufficiently large $k$, this $y$ belongs to $K_{2R}$ whenever
$x\in K_R$. Indeed the compact set $K_R$ lies in the interior of
$K_{2R}$, and $\|y-x\|\le\sigma_k\to0$. It is also generic by the
strict coordinate separation just proved. The free function $Q$ is
uniformly continuous on $K_{2R}$. If $\omega_Q(t)$ is its modulus of
continuity there, the three-term triangle inequality gives
\begin{equation}\label{eq:three-clear}
 \sup_{x\in K_R\cap\Omega_{\sigma_k}}
 |U_{k,Y}(x)-Q(x)|
 \le\alpha_k:=16m^2B_k^m h_k/\sigma_k+e_k+\omega_Q(\sigma_k)
 \longrightarrow0.
\end{equation}
This convergence is uniform over the limiting tuples $Y$ under consideration.

There is an explicit linear bound for the volume of the omitted strips:
\begin{equation}\label{eq:three-strip}
 |K_R\setminus\Omega_\sigma|\le C_{m,R}\sigma,
 \qquad C_{m,R}=64m(m-1)(2R)^{4m-1}.
\end{equation}
For a fixed ordered pair $i\ne j$ and coordinate $\nu$, make the
linear change replacing $x_{j,\nu}$ by
$x_{j,\nu}-x_{i,\nu}$ and leaving the other coordinates unchanged.
Its determinant is one. The restricted coordinate has interval length
$16\sigma$ and each of the remaining $4m-1$ coordinates lies in an
interval of length $2R$. This bounds the volume of that strip intersected
with $K_R$ by $16\sigma(2R)^{4m-1}$. There are $4m(m-1)$ choices.
Subadditivity proves \eqref{eq:three-strip}, including intersecting strips.
For $m=1$ the union is empty and the constant is zero.

Let $f$ be a smooth test supported in $K_R$, and put
$C_Q=\sup_{K_R}|Q|<\infty$. Split its integral into the clear region
and its complement. Equations \eqref{eq:three-bulk}--\eqref{eq:three-strip}
give the quantitative bound
\begin{equation}\label{eq:three-integral}
 \left|\int f(x)\{U_{k,Y}(x)-Q(x)\}\,dx\right|
 \le\|f\|_\infty\left\{|K_R|\alpha_k+
                  C_{m,R}(A_k^m+C_Q)\sigma_k\right\}.
\end{equation}
The right side tends to zero. Both integrands are continuous on a
neighborhood of the support, so these are ordinary integrals; no extension
across coincident configurations is involved.

\subsection{Returning to cutoff-dependent insertions}

It is necessary to apply this estimate to the actual insertion maps, whose
values depend on the cutoff. Fix $k$. Their eventual boundedness gives a
compact closure in their finite-dimensional operator space. We use the
ordered point and smeared limits of the same table $S$; these limits also
hold on subfilters. If insertion maps converge to a map $T$ along a
subfilter, multilinearity in a fixed finite basis proves that their point
and smeared correlations converge to those with the limiting labels.
For example, expanding an $m$-tuple gives finitely many terms
\[
 \sum_{a_1,\ldots,a_m}
     t_{1,a_1}(N,p)\cdots t_{m,a_m}(N,p)
     U_{N,p}(e_{a_1},\ldots,e_{a_m}).
\]
Every scalar coefficient converges and every fixed basis correlation has
the specified limit. Finite sums and products therefore give the claimed
convergence for varying insertions. For the smeared expression, use the
fixed dilated test $r_k^{-4m}f(\,\cdot\,/r_k)$; a change of variables identifies
its limiting integral with $\int f(x)U_{k,Y}(x)\,dx$.

Denote the right side of \eqref{eq:three-integral} by $E_k(f)$. We claim
that the cutoff smeared comparison error is eventually at most
$E_k(f)+\eta$, for every $\eta>0$. If not, the set where it exceeds this
number is frequent in $\mathcal F$. Restricting to that set gives a
nontrivial filter. Compactness of the insertion maps supplies a further
convergent subfilter. Along it the smeared error converges, by the preceding
finite expansion, to a number at most $E_k(f)$. But its values are all
larger than $E_k(f)+\eta$, so their limit is at least that latter number.
This is a contradiction. Consequently the actual cutoff errors have the
same asymptotic upper bound, and that bound tends to zero as $k\to\infty$.
This compactness argument is performed separately at each fixed $k$ and
does not replace the original ordered extraction by a new continuum table.

\subsection{The three scale inequalities}

We describe precisely the estimates needed from the marked construction.
Fix an upper arity $n\ge2$, a maximum grade $D$, and a word of length
$m\le n$ in grades at most $D$. Let $H_k>1$ be the diverging comparison
monitor. Choose three exponents $a,b,e>0$ and a strip exponent $c<0$,
and set
\[
 r_k=a_kH_k^a,\qquad h_k=H_k^{-a},\qquad
 \sigma_k=H_k^c.
\]
At a fixed comparison depth, all three panels must use the same calibrated
readout map $T$. The lift at exponent $a$ and the lift at exponent $b$
then differ by the grade weight $H_k^{(a-b)j}$ on grade $j$.
With a common finite basis, the sum of the absolute expansion coefficients
of a fixed label of grade at most $D$ is bounded by a constant times
$H_k^{(a-b)D}$. The constant depends on the label, not on $k$.
This uses the same map $T$ on the two sides; two unrelated calibrations
would not give the stated identity.

The marked comparison at the $a$-panel supplies
\eqref{eq:three-grid} with $e_k=C\varepsilon_k\to0$.
The centered one-point bounds at the $b$-panel, expanded in that common
basis and in the finitely many signed and conjugated orientations, supply
the bulk coefficient bound with
\[
          A_k=C_bH_k^{(a-b)D}.
\]
The analogous centered difference bound at the $e$-panel supplies
\[
          B_k=C_eH_k^{(a-e)D}.
\]
Here the product and difference estimates are the centered cell estimates
of Section~\ref{ex-centered-estimates}. Their use requires actual cell
widths, not just these coefficient prices. If $w_b,w_e$ are the output
levels of the latter two panels, the schedule estimates are
\begin{gather}
 a_{w_b-1}/a_k\le H_k^{b+6},\qquad
          6m a_{w_b}<a_{w_b-1},\label{eq:three-width-b}\\
 a_{w_e}/a_k\le H_k^{e+5}.\label{eq:three-width-e}
\end{gather}
The first gives a bulk cell strictly smaller than $r_k/(2R)$ once
$b+6<a$; the second gives a difference cell smaller than
$r_k\sigma_k$ once $e+5<a+c$. Applying the centered bounds and their
continuum limits then gives \eqref{eq:three-bulk} and the directional
bound $B_k^m/\sigma_k$. Summing over the $4m$ coordinate directions
gives exactly \eqref{eq:three-slope}.

Here is the elementary rounding calculation underlying the width expenses
$5$ and $6$. On the relevant constant-block interval let its block factor
be $L\le H_k$ and let $V$ be its cover parameter, with $V^{10}\le H_k$.
The panel height is the least $h$ with $L^h\ge H_k^{b+2}$, and its
surplus is the least $s$ with $L^s\ge V^{10}$. Minimality gives
$L^h\le L H_k^{b+2}$ and $L^s\le L V^{10}$. Thus
\[
 a_{w_b}/a_k=L^{h+s}\le H_k^{b+5}.
\]
One earlier block costs at most another $H_k$, giving the first inequality
in \eqref{eq:three-width-b}. The block factors at these output levels
tend to infinity; hence they eventually exceed $6m$, proving its second
inequality. The identical calculation with $e$ gives
\eqref{eq:three-width-e}. The allowance of each panel must place all
these levels in the same constant-block interval.

It remains to choose the exponents so that both integral errors decay.
The exact sufficient inequalities are
\begin{equation}\label{eq:three-checks}
 \begin{split}
 c&<0,\qquad b+6<a,\qquad e+5<a+c,\\
 (a-b)Dn+c&<0,\qquad (a-e)Dn-a-c<0.
 \end{split}
\end{equation}
They imply $-a<c$ and hence $4m h_k\le\sigma_k$ eventually.
They also give
\[
 A_k^m\sigma_k=C_b^mH_k^{(a-b)Dm+c}\to0,\qquad
 B_k^m h_k/\sigma_k=C_e^mH_k^{(a-e)Dm-a-c}\to0.
\]
Indeed $a>b,e$, so replacing $m$ by $n$ increases the exponents.
For a concrete choice, write $q=Dn$, $x=7q+1$, $y=6+x$, and take
\[
 a=10+D+x+yq+y+7,\qquad b=a-7,\qquad e=a-y,\qquad c=-x.
\]
Then $b+6=a-1$, $e+5=a-x-1$, and the last two exponents in
\eqref{eq:three-checks} are respectively $-1$ and $-17-D-y$.
Also $a,b,e>D+6$, so the additional lower bounds on panel exponents
are respected. These choices depend only on the finite grade and arity
bounds. They make no assertion of a uniform estimate over all words.

Finally, if these common-panel data are available at arbitrarily large
comparison depths, choose an increasing sequence of those depths.
All the eventually valid numerical inequalities remain valid along it.
There are only finitely many basis words under fixed grade and arity
bounds; take the maximum of their finitely many thresholds. The preceding
proof applies to each word on this one sequence, and multilinearity extends
it to arbitrary labels within these bounds. The empty word is the constant
one and needs no strip estimate. Unit, conjugation and signed-symmetry
properties of the insertion maps pass through this construction whenever
they hold eventually at the finite cutoffs. Producing the common marked
panel data on the selected Wilson history is a separate upstream task;
the argument here proves their exact weak-convergence consequence.
